\cotefile{18}
\scanpage{48}
% archive sheet number: 47
\mspage{2 - 12 -}\hfill n\textsuperscript{o} 392

c) Le corps résiduel de $v$ s'identifie à $A/\unsure{\underset{=}{p}|\underline{p}|p}$, et celui de $v_i$ à $B/\underset{=}{p}_i$, (§ 1, n\textsuperscript{o} 4, cor. 1 de la prop. 6).

d) Soit $a$, (resp. $a_i$), une uniformisante pour $v$, (resp. $v_i$). On a
\[
\begin{aligned}
aB_{\underset{=}{p}_i}&=aA_{\underset{=}{p}}B_{\underset{=}{p}_i}
=\underset{=}{p}A_{\underset{=}{p}}B_{\underset{=}{p}_i}
=\underset{=}{p}B.B_{\underset{=}{p}_i}
=(\handwritten{\prod}_{j}\underset{=}{p}_j^{e(j)})B_{\underset{=}{p}_i}
=\handwritten{\prod}_{j}(\underset{=}{p}_j B_{\underset{=}{p}_i})^{e(j)}
=(\underset{=}{p}_i B_{\underset{=}{p}_i})^{e(i)}=\\
&=a_i^{e(i)}B_{\underset{=}{p}_i}\quad\text{puisque }\underset{=}{p}_jB_{\underset{=}{p}_i}=B_{\underset{=}{p}_i}\quad\text{pour }j\handwritten{\ne}i\quad\text{; d'où d), puisque}\\
&e(v_i/v)=v_i(a)\,.
\end{aligned}
\]

\begin{center}-:-:-:-\end{center}
